\chapter{表示、计算与不准确的模型}
\label{ch:representation}
\chapterguide{保存区别、低成本读出区别，以及正确描述区别，是三个需要分别证明的条件。}

\section{组织不能把未知答案藏进新操作}
假设将$n$步协议的最终输出重新命名为一次宏操作。它形成一个新概率核，但执行这个核仍需原来的$n$次调用。将宏操作当作单位成本，会人为改变能力比较。

对任意两情境，宏输出的KL不大于完整记录的KL，而完整记录KL等于第\ref{ch:lower}章的调用账本。因此后处理可以降低计算或通信负担，不能绕过原有证据下界。若某种“检查”真的提供了此前没有的观察分布，它应作为新基本能力计费。

\section{同样的信息可以需要不同读出成本}
考虑已经观察到的$k\ge2$个独立公平bit，目标是它们的异或。一个只能读取一个原坐标的读出器，错误率为$1/2$；先计算异或并把结果放入一个新片段，同一类单坐标读出器便能零错误读取目标。

原数据与新片段都含有一bit目标信息。变化发生在允许的读出计算，而异或计算已经由准备片段的过程付出$O(k)$的费用。可利用信息理论明确将预测器的能力纳入度量\citep{xu2020}。这里的诊断定理首先给出统计限制，实现时再列出取得相应解码器所需成本。

第\ref{ch:calibration}章的消元构造具有不同作用：它在显式实验矩阵上找到目标的表达，再用有噪观察执行表达。由此避免直接枚举$2^r$个情境。表示的收益通过可执行算法体现，并非仅给原问题换一个符号。

\section{概率表估计错误会消耗多少余量}
实际解码使用估计表$\widehat K$，真实表仍为$K$。设所有概率严格正，并有统一对数误差
\[
 |\ln\widehat K_a(v\mid\vartheta)-\ln K_a(v\mid\vartheta)|\le\varepsilon.
\]
分配$w$在真实表上具有最差覆盖率$u$。估计最大似然输出按$\widehat L$选择情境。

\begin{theorem}[模型误差的可靠性代价]
\label{thm:robust}
上述组织满足
\[
 \max_\theta\PP_\theta(g(\widehat\theta)\ne g(\theta))
 \le\min\{1,me^{n\varepsilon}(1-u)^n\}
 \le\min\{1,me^{-n(u-\varepsilon)}\}.
\]
\end{theorem}
\begin{proof}
每个候选的累计对数似然误差至多$n\varepsilon$。若错误候选在估计表上胜出，则真实对数似然差$\ln L_n(\theta')-\ln L_n(\theta)$至少为$-2n\varepsilon$。该事件的概率至多$e^{n\varepsilon}$乘两种真实记录分布的亲和度。再对至多$m$个候选求和。
\end{proof}

若$u>\varepsilon$，这个保证仍给出指数压低误差的构造。若$u\le\varepsilon$，所写上界不再保证收敛；这不等价于已经证明所有算法失败。统一对数误差需要控制小概率项，有限样本平均拟合误差无法自动替代它。

同一证明也给通信量化界：若接收者累计的每个对数似然分数误差至多$E$，最大似然风险至多$me^E(1-u)^n$。这允许根据可用诊断余量选择摘要精度。仅仅传递浮点分数而不规定误差，会使所谓“充分消息”失去保证。

\section{模型失配的更深层问题}
上述定理要求真实情境已在候选集合中。若全部候选共享错误解释，增加观察会可靠地选择一个失配模型，不能得到关于真实目标的保证。

对一组可能概率表，可先检查它们是否包含两个目标不同但全部观察分布相同的实例。一旦存在这样的对，定理\ref{thm:blind}直接给出稳健推断的障碍。不存在这类对也未必意味着容易学习：还要解决实验设计、有限样本辨认与模型表示成本。

\begin{exercise}每个成员提供的各候选对数似然分数误差至多$e_j$。给出汇总后风险上界，并说明误差预算如何相加。\end{exercise}
\begin{exercise}如果原始片段的内容只经过可逆重编码，为什么信息区别不变，而实际解码时间仍可能改变？\end{exercise}
